\chapter{职业性致癌因素与职业性肿瘤}
\setcounter{chapter}{6}
\chapterintro{
\textbf{【教学大纲要求】}

\imp{掌握：}职业性肿瘤、致癌因素的概念；人类致癌物、可疑致癌物和潜在致癌物；我国法定的职业肿瘤；职业性致癌因素的分类；职业性致癌因素的识别与判断；职业性肿瘤的特征\keyword{*}（病因、潜伏期、阈值、好发部位、病理类型、影响因素）；常见的职业肿瘤（职业性呼吸道癌、皮肤癌、膀胱癌、白血病、肝血管肉瘤）。

\imp{熟悉：}职业性肿瘤的诊断及处理原则；职业性肿瘤的预防原则（法律法规、监测体系、危险管理、控制管理、健康教育、化学预防）。

（注：\keyword{*} 标示教学难点）
}

% ----- 6.1 概述 -----
\section{职业性致癌因素}

\begin{defbox}
\textbf{职业性致癌因素}：在劳动过程中，可能引起劳动者发生职业性肿瘤的各种因素。包括化学因素（最常见）、物理因素和生物因素。

\textbf{分类：}
\begin{enumerate}[leftmargin=*]
    \item \imp{确认致癌物}：流行病学调查和动物实验均有充分证据证明对人致癌
    \item \imp{可疑致癌物}：动物实验证据充分，对人体致癌证据有限（IARC 2A类）
    \item \imp{潜在致癌物}：动物实验已获阳性结果但人群流行病学调查尚无明确证据（IARC 2B类）
\end{enumerate}
\end{defbox}

% ----- 6.2 职业性肿瘤的特点 -----
\section{职业性肿瘤的特点*}

\begin{keybox}
\textbf{1. 潜伏期长}：从首次接触到临床发病通常需要10--30年或更长（平均约20年）。因此多数职业性肿瘤发生在接触者的中老年期。

\textbf{2. 阈值问题}：对于遗传毒性致癌物，目前多认为不存在绝对安全的阈值。但对于非遗传毒性致癌物，可能存在阈值。

\textbf{3. 剂量-反应关系}：多数职业性肿瘤表现出明确的剂量-反应关系——接触剂量越大、时间越长，发病风险越高。

\textbf{4. 好发部位}：具有较固定的好发部位——多发生于致癌物接触最直接、最强烈的部位：
\begin{itemize}[leftmargin=*]
    \item 呼吸道（吸入性致癌物）：石棉→肺癌、间皮瘤
    \item 皮肤（皮肤接触）：煤焦油→皮肤癌
    \item 膀胱（经尿排泄）：联苯胺→膀胱癌
    \item 肝脏（代谢活化部位）：氯乙烯→肝血管肉瘤
    \item 造血系统：苯→白血病
\end{itemize}

\textbf{5. 影响因素*：}
\begin{itemize}[leftmargin=*]
    \item 吸烟（协同作用）：石棉暴露+吸烟→肺癌风险增加数十倍
    \item 年龄（初次接触年龄越轻，风险越高）
    \item 个体遗传易感性（代谢酶基因多态性、DNA修复能力差异）
\end{itemize}

\textbf{6. 病理类型}：某些职业性肿瘤具有特定的病理类型。例如氯乙烯所致肝血管肉瘤、联苯胺所致膀胱癌（移行上皮癌为主）。
\end{keybox}

% ----- 6.3 我国法定职业性肿瘤 -----
\section{我国法定职业性肿瘤（11种）}

\begin{keybox}
\begin{table}[H]
\centering
\caption{我国法定职业性肿瘤（11种）}
\begin{tabularx}{\textwidth}{ll}
\hline
\textbf{职业性肿瘤} & \textbf{接触因素} \\
\hline
1. 石棉所致肺癌、间皮瘤 & 石棉 \\
2. 联苯胺所致膀胱癌 & 联苯胺 \\
3. 苯所致白血病 & 苯 \\
4. 氯甲醚、双氯甲醚所致肺癌 & 氯甲醚、双氯甲醚 \\
5. 砷及其化合物所致肺癌、皮肤癌 & 砷及其无机化合物 \\
6. 氯乙烯所致肝血管肉瘤 & 氯乙烯 \\
7. 焦炉逸散物所致肺癌 & 焦炉逸散物（含PAHs）\\
8. 六价铬化合物所致肺癌 & 六价铬化合物 \\
9. 毛沸石所致肺癌、间皮瘤 & 毛沸石 \\
10. 煤焦油、煤焦油沥青、石油沥青所致皮肤癌 & 煤焦油、沥青 \\
11. $\beta$-萘胺所致膀胱癌 & $\beta$-萘胺 \\
\hline
\end{tabularx}
\end{table}
\end{keybox}

% ----- 6.4 职业性肿瘤的识别和预防 -----
\section{职业性致癌因素的识别与判定}

\begin{enumerate}[leftmargin=*]
    \item \imp{临床观察}：临床医生通过病例报告发现可疑致癌因素
    \item \imp{流行病学调查}：队列研究、病例对照研究等——\textbf{是识别和确认职业性致癌因素中\keyword{最重要的一步}，最具说服力}
    \item \imp{动物实验}：致癌性生物鉴定（长期致癌试验）
    \item \imp{短期试验}：Ames试验（致突变试验）等
    \item \imp{分子流行病学}：生物标志物（DNA加合物、基因突变谱等）
\end{enumerate}

\begin{tipbox}
\textbf{职业性致癌因素的判定标准（WHO/IARC分类）：}
\begin{itemize}[leftmargin=*]
    \item \textbf{IARC 1类（确定致癌物）}：对人类致癌的充分证据（如石棉、苯、六价铬、氯乙烯等）
    \item \textbf{IARC 2A类（很可能致癌物）}：有限的人群证据+充分的动物实验证据
    \item \textbf{IARC 2B类（可能致癌物）}：有限的人群证据+不充分的动物证据
    \item \textbf{IARC 3类}：无法分类（证据不足）
\end{itemize}
\textbf{注意：}流行病学调查证据是判定确认致癌物的根本依据。仅凭动物实验阳性不足以将某因素列为确认致癌物。
\end{tipbox}

\begin{exambox}
\textbf{工程措施与管理措施的区分：}
\begin{itemize}[leftmargin=*]
    \item \textbf{工程措施}：通过技术手段从根本上减少或消除危害，如密闭化生产自动化、通风排毒、湿式作业、隔声吸声、密闭尘源、无毒代有毒等
    \item \textbf{管理措施/组织措施}：通过制度和管理手段控制危害，如合理安排作息时间、轮换作业、安全操作规程、健康监护制度、安全教育培训等
    \item 工程措施是\textbf{首选}的预防策略（从源头消除危害），管理措施和个人防护为辅助手段
\end{itemize}
\end{exambox}

\section{职业性肿瘤的预防原则}

\begin{enumerate}[leftmargin=*]
    \item 严格执行有关职业卫生标准和法律法规
    \item 建立和完善有效的职业性肿瘤监测体系
    \item 落实预防致癌危害的管理办法
    \item 加强对职业性致癌因素的控制和管理
    \item 加强健康教育提高自我防护能力
    \item 科学预防：识别高危人群、定期医学监护
\end{enumerate}

% ----- 6.5 复习题 -----
\section{复习题}

\begin{enumerate}[leftmargin=*, itemsep=8pt]
    \item \textbf{【名词解释】}职业性致癌因素；确认致癌物；可疑致癌物
    \item \textbf{【简答题】}简述职业性肿瘤的六大特点。
    \item \textbf{【简答题】}列举至少5种我国法定职业性肿瘤及其接触因素。
    \item \textbf{【简答题】}简述职业性致癌因素的识别与判定方法。
    \item \textbf{【简答题】}简述职业性肿瘤的预防原则。
\end{enumerate}

\subsection*{参考答案要点}

\begin{enumerate}[leftmargin=*]
    \item \textbf{职业性致癌因素}：劳动过程中可能引起职业性肿瘤的各种因素（化学、物理、生物）。\textbf{确认致癌物}：流行病学和动物实验均充分证明对人致癌。\textbf{可疑致癌物}：动物证据充分但人体证据有限（2A类）。
    \item 六大特点：潜伏期长（10--30年）；阈值问题（遗传毒性致癌物无安全阈值）；明确的剂量-反应关系；好发部位较固定（接触最直接的部位）；吸烟和个体易感性为重要影响因素；某些肿瘤有特定病理类型。
    \item 石棉→肺癌/间皮瘤；联苯胺→膀胱癌；苯→白血病；氯甲醚→肺癌；砷→肺癌/皮肤癌；氯乙烯→肝血管肉瘤；焦炉逸散物→肺癌；六价铬→肺癌；煤焦油沥青→皮肤癌。
    \item 临床观察→流行病学调查（最具说服力）→动物实验（长期致癌试验）→短期试验（Ames试验等）→分子流行病学（生物标志物）。
    \item 执行法律法规标准；建立肿瘤监测体系；落实预防致癌危害管理办法；加强致癌因素控制管理；加强健康教育；科学预防（识别高危人群、定期监护）。
\end{enumerate}

\newpage

% =====================================================================
%  CHAPTER 7: 生物性有害因素所致职业性损害
% =====================================================================
\newpage